% source: J. S. Milne, "Abelian Varieties" (course notes), v2.00, 2008, www.jmilne.org
% pdf_page: 0021
% doc_page: 15 (Chapter I, end of Section 2 "Abelian Varieties over the Complex Numbers" — dual torus, Weil pairing, polarizations, Notes; start of Section 3 "Rational Maps Into Abelian Varieties")
% retrieved: 2026-07-17
% transcribed_by: claude-fable-5 (vision, rendered at 200dpi via pdftoppm)

% [running head: 3. RATIONAL MAPS INTO ABELIAN VARIETIES, page 15]

\DeclareMathOperator{\End}{End}     % End(E), End(X), endomorphism algebra
\DeclareMathOperator{\Image}{Im}    % Im f(L), imaginary part

% [continuation of Section 2 (complex tori) from the previous page]

Let $X = V/L$. Then
\[
  V^* = \{ f \colon V \to \mathbb{C} \mid f(\alpha v) = \bar{\alpha} f(v),
    \quad f(v + v') = f(v) + f(v') \}
\]
is a complex vector space of the same dimension as $V$. Define
\[
  L^* = \{ f \in V^* \mid \Image f(L) \subset \mathbb{Z} \}.
\]
Then $L^*$ is a lattice in $V^*$, and
$X^{\vee} \overset{\text{df}}{=} V^*/L^*$ is a polarizable complex torus,
called the \textbf{dual torus}.

\paragraph{Exercise 2.11.}
If $X = V/L$, then $X_m$, the subgroup of $X$ of elements killed by $m$, is
$m^{-1}L/L$. Show that there is a canonical pairing
\[
  X_m \times (X^{\vee})_m \to \mathbb{Z}/m\mathbb{Z}.
\]
This is the \textbf{Weil pairing}.

A Riemann form on $E$ on $X$ defines a homomorphism
% [sic: printed "A Riemann form on E on X"; presumably "A Riemann form E on X" is intended]
$\lambda_E \colon X \to X^{\vee}$ as follows: let $H$ be the associated
Hermitian form, and let $\lambda_E$ be the map defined by
\[
  v \mapsto H(v, \cdot) \colon V \to V^*.
\]
Then $\lambda_E$ is an isogeny, and we call such a map $\lambda_E$ a
\textbf{polarization}. The \textbf{degree} of the polarization is the order
of the kernel. The polarization is said to be \textbf{principal} if it is of
degree 1.

\paragraph{Exercise 2.12.}
Show that every polarizable tori is isogenous to a principally polarized
torus.
% [sic: printed "every polarizable tori"; presumably "every polarizable torus" is intended]

A polarizable complex torus is \textbf{simple} if it does not contain a
nonzero proper polarizable subtorus $X'$.

\paragraph{Exercise 2.13.}
Show that every polarizable torus is isogenous to a direct sum of simple
polarizable tori.

Let $E$ be an elliptic curve over $\mathbb{Q}$. Then
$\End(E) \otimes \mathbb{Q}$ is either $\mathbb{Q}$ or a quadratic imaginary
extension of $\mathbb{Q}$. For a simple polarizable torus,
$D = \End(X) \otimes \mathbb{Q}$ is a division algebra over a field and the
polarization defines a positive involution ${}^{\dagger}$ on $D$. The pairs
$(D, {}^{\dagger})$ that arise from a simple abelian variety have been
classified (A.A. Albert).

\subsection*{Notes}

There is a concise treatment of complex abelian varieties in Chapter I of
Mumford 1970, and a more leisurely account in Murty 1993. The classic
account is Siegel 1948. Siegel first develops the theory of complex
functions in several variables. See also his books, Topics in Complex
Function Theory. There is a very complete modern account in Birkenhake and
Lange 2004
% [sic: the paragraph ends "Birkenhake and Lange 2004" with no final period in the printed text]

\section*{3 \quad Rational Maps Into Abelian Varieties}

Throughout this section, all varieties will be irreducible.
